% ============ 第1章 §1.1 准备(原书 pp.1-9 = PDF pp.13-21) ============
\section{准备}\label{sec:1.1}
\index{Lebesgue 收敛定理}
\index{积分号下的微分定理}
\index{Fubini 定理}
\index{Hölder 不等式}
\index{Cauchy--Schwarz 不等式}
\index{Young 不等式}
\index{Schwartz 函数}
\index{速降函数}


首先, 列举经常使用的记号.以下, $A := B$ 表示将 $B$ 定义为 $A$.

$\mathbb{R} := (-\infty, \infty)$, $\mathbb{R}_{+} := (0, \infty)$, $\mathbb{C}$:全体复数

\begin{equation*}
\mathbb{N} := \{1, 2, 3, \dots \}, \quad
\mathbb{N}_{0} := \{0, 1, 2, 3, \dots \}, \quad
\mathbb{Z} := \{\dots, -2, -1, 0, 1, 2, \dots \}
\end{equation*}

$C(\mathbb{R})$:$\mathbb{R}$ 上连续函数的全体

$C_0(\mathbb{R})$:$\mathbb{R}$ 上满足 $\lim_{|x|\to\infty}f(x) = 0$ 的连续函数 $f$ 的全体

$C_{c}(\mathbb{R})$:$\mathbb{R}$ 上支集为有界闭集的连续函数的全体

$C^{k}(\mathbb{R})$:$\mathbb{R}$ 上 $k$ 次连续可微函数的全体

$L^{\infty}(\mathbb{R})$:$\mathbb{R}$ 上本性有界函数的全体

\begin{equation*}
\| f \|_{L^{\infty}(\mathbb{R})} = \| f \|_{L^{\infty}} = \| f \|_{\infty} := \operatorname{ess\,sup}_{x \in \mathbb{R}} |f(x)|
\end{equation*}

$L^{p}(\mathbb{R})$:$\mathbb{R}$ 上 $p$ 次可积函数的全体

\begin{equation*}
\| f \|_{L^{p}(\mathbb{R})} = \| f \|_{L^{p}} = \| f \|_{p} := \left(\int_{\mathbb{R}} |f(x)|^{p}\dd x\right)^{1/p}
\end{equation*}

$\lim_{j\to\infty}\|f_{j}-f\|_{L^{p}} = 0$ 记为 $f_{j}\to f$ in $L^{p}$($j\to\infty$).

对几乎所有的 $x \in \mathbb{R}$ 成立 $f(x) = g(x)$ 记为
$f(x) = g(x)$(a.e. $x\in\mathbb{R}$).

$L_{\mathrm{loc}}^{p}(\mathbb{R})$:在 $\mathbb{R}$ 的任意有界区间上 $p$ 次可积的函数的全体

对 $\mathbb{R}$ 上的可测集 $E$, 用 $|E|$ 表示 $E$ 的 Lebesgue 测度.

\begin{definition}\label{def:1.1.1}
$f \in \mathcal{S}(\mathbb{R})$, 或者说 $f$ 是 Schwartz 函数, 或者说 $f$ 是速降函数, 是指 $f \in C^{\infty}(\mathbb{R})$, 并且对任意 $k, \ell \in \mathbb{N}_0$, 存在正数 $C_{k,\ell}$ 使得
\begin{equation*}
\left| (1 + |x|)^{k} f^{(\ell)}(x) \right| \leq C_{k,\ell} \quad (x \in \mathbb{R}).
\end{equation*}
\end{definition}

接下来, 叙述 Lebesgue 积分中的基本事实.

\begin{theorem}[Lebesgue 收敛定理]\label{thm:1.1.2}
$f_{j}(x)$($j = 1, 2, \ldots$)是可测函数, $\varphi(x) \in L^{1}(\mathbb{R})$, 且
\begin{align*}
|f_{j}(x)| \leq \varphi(x) & (\text{a.e. } x \in \mathbb{R},\ j = 1, 2, \ldots) \\
\lim_{j \to \infty} f_{j}(x) = f(x) & (\text{a.e. } x \in \mathbb{R})
\end{align*}
时, 
\begin{equation*}
\lim_{j \rightarrow \infty} \int_{-\infty}^{\infty} f_{j}(x)\dd x = \int_{-\infty}^{\infty} f(x)\dd x.
\end{equation*}
\end{theorem}

\begin{corollary}\label{cor:1.1.3}
设 $1 \leq p < \infty$, $f_{j}(x)$($j = 1, 2, \ldots$)是可测函数, $\varphi(x) \in L^{p}(\mathbb{R})$, 且
\begin{align*}
|f_{j}(x)| \leq \varphi(x) & (\text{a.e. } x \in \mathbb{R},\ j = 1, 2, \ldots) \\
\lim_{j \to \infty} f_{j}(x) = f(x) & (\text{a.e. } x \in \mathbb{R})
\end{align*}
时, 
\begin{align*}
\lim_{j \to \infty} \left(\int_{-\infty}^{\infty} |f_{j}(x) - f(x)|^{p}\dd x\right)^{1/p} &= 0, \\
\lim_{j \to \infty} \left(\int_{-\infty}^{\infty} |f_{j}(x)|^{p}\dd x\right)^{1/p} &= \left(\int_{-\infty}^{\infty} |f(x)|^{p}\dd x\right)^{1/p}.
\end{align*}
\end{corollary}

\begin{theorem}[积分号下的微分定理]\label{thm:1.1.4}
设 $f(x,t)$ 关于 $t$ 可偏导, 对每个 $t$, $f(x,t)$ 与 $f_{t}(x,t)$ 关于 $x$ 可积, 并且存在可积函数 $\varphi(x)$ 使得 $|f_{t}(x,t)| \leq \varphi(x)$.于是 $\int_{-\infty}^{\infty} f(x,t)\dd x$ 关于 $t$ 可微, 且
\begin{equation*}
\frac{\dd}{\dd t}\left(\int_{-\infty}^{\infty} f(x,t)\dd x\right) = \int_{-\infty}^{\infty} f_{t}(x,t)\dd x.
\end{equation*}
这里, 若对几乎所有的 $x$, $f_{t}(x,t)$ 关于 $t$ 连续, 则微分得到的函数也连续.
\end{theorem}

\begin{proof}
令 $F(t) = \int_{-\infty}^{\infty} f(x,t)\dd x$, 则
\begin{equation*}
\frac{F(t+h) - F(t)}{h} = \int_{-\infty}^{\infty} \frac{f(x,t+h) - f(x,t)}{h}\dd x.
\end{equation*}
由中值定理, 
\begin{equation*}
\left| \frac{f(x,t+h) - f(x,t)}{h} \right| \leq \varphi(x),
\end{equation*}
又由假设
\begin{equation*}
\lim_{h \rightarrow 0} \frac{f(x,t+h) - f(x,t)}{h} = f_{t}(x,t),
\end{equation*}
故由 Lebesgue 收敛定理
\begin{equation*}
F^{\prime}(t) = \lim_{h \rightarrow 0} \frac{F(t+h) - F(t)}{h} = \int_{-\infty}^{\infty} f_{t}(x,t)\dd x.
\end{equation*}
最后一步也是 Lebesgue 收敛定理的应用.
\end{proof}

\begin{theorem}[Fubini 定理]\label{thm:1.1.5}
设 $f(x,y)$ 关于 $(x,y)$ 可测.若
\begin{align*}
&\iint_{\mathbb{R}^{2}} |f(x,y)|\dd x\dd y, \quad \int_{-\infty}^{\infty}\left(\int_{-\infty}^{\infty} |f(x,y)|\dd x\right)\dd y, \\
&\int_{-\infty}^{\infty}\left(\int_{-\infty}^{\infty} |f(x,y)|\dd y\right)\dd x
\end{align*}
三者之一有限, 则其余两个也有限, 并且成立如下结论:

(i) 对几乎所有的 $y \in \mathbb{R}$, $\int_{-\infty}^{\infty} f(x,y)\dd x$ 存在且作为 $y$ 的函数可测.

(ii) 对几乎所有的 $x \in \mathbb{R}$, $\int_{-\infty}^{\infty} f(x,y)\dd y$ 存在且作为 $x$ 的函数可测.

(iii)
\begin{align*}
\iint_{\mathbb{R}^{2}} f(x,y)\dd x\dd y
& = \int_{-\infty}^{\infty}\left(\int_{-\infty}^{\infty} f(x,y)\dd x\right)\dd y \\
& = \int_{-\infty}^{\infty}\left(\int_{-\infty}^{\infty} f(x,y)\dd y\right)\dd x.
\end{align*}
\end{theorem}

\begin{corollary}\label{cor:1.1.6}
设 $0 < p < \infty$.则
\begin{equation*}
\int_{-\infty}^{\infty} |f(x)|^{p}\dd x = p\int_{0}^{\infty} |\{x;\ |f(x)| > t\}|\,t^{p-1}\dd t.
\end{equation*}
\end{corollary}

\begin{theorem}[Hölder 不等式]\label{thm:1.1.7}
设 $1 \leq p \leq \infty$, $1/p + 1/p' = 1$.若 $f(x) \in L^{p}(\mathbb{R})$, $g(x) \in L^{p'}(\mathbb{R})$, 则 $f(x)g(x)$ 可积, 且
\begin{equation*}
\left| \int_{-\infty}^{\infty} f(x)g(x)\dd x \right| \leq \| f \|_{L^{p}} \| g \|_{L^{p'}}.
\end{equation*}
\end{theorem}

Hölder 不等式在 $p = 2$ 时称为 Cauchy--Schwarz 不等式.此后也常以此名使用.

\begin{theorem}[Young 不等式的一个特殊情形]\label{thm:1.1.8}
设 $1 \leq p \leq \infty$.若 $f(x) \in L^{1}(\mathbb{R})$, $g(x) \in L^{p}(\mathbb{R})$, 则对几乎所有的 $x$, $\int f(x-y)g(y)\dd y$ 存在, 且下式成立:
\begin{equation*}
\left(\int_{-\infty}^{\infty}\left| \int_{-\infty}^{\infty} f(x-y)g(y)\dd y \right|^{p}\dd x\right)^{1/p} \leq \| f \|_{L^{1}} \| g \|_{L^{p}}.
\end{equation*}
\end{theorem}

\begin{proof}
\begin{align*}
& \left| \int_{-\infty}^{\infty} f(x-y)g(y)\dd y \right| \\
& \qquad \leq \int_{-\infty}^{\infty} |f(x-y)g(y)|\dd y
\end{align*}
\begin{align*}
&\leq \left(\int_{-\infty}^{\infty} |f(x-y)|\dd y\right)^{1/p'}\left(\int_{-\infty}^{\infty} |f(x-y)||g(y)|^{p}\dd y\right)^{1/p} \\
&= (\| f \|_{L^{1}})^{1/p'}\left(\int_{-\infty}^{\infty} |f(x-y)||g(y)|^{p}\dd y\right)^{1/p}.
\end{align*}
又由 Fubini 定理
\begin{align*}
\int_{-\infty}^{\infty} & \left(\int_{-\infty}^{\infty} |f(x-y)||g(y)|^{p}\dd y\right)\dd x
= \int_{-\infty}^{\infty}\left(\int_{-\infty}^{\infty} |f(x-y)||g(y)|^{p}\dd x\right)\dd y \\
& = \| f \|_{L^{1}} \int_{-\infty}^{\infty} |g(y)|^{p}\dd y < \infty.
\end{align*}
从而, 由 Fubini 定理, 对几乎所有的 $x$, $\int_{-\infty}^{\infty} f(x-y)g(y)\dd y$ 存在.由上面的两个不等式即得所证不等式.
\end{proof}

\begin{theorem}\label{thm:1.1.9}
设 $K(t,x)$ 满足如下三条:
\begin{enumerate}[(i)]
\item $\int_{-\infty}^{\infty}|K(t,x)|\dd x \leq A \quad (t > 0)$, 
\item $\int_{-\infty}^{\infty} K(t,x)\dd x = 1 \quad (t > 0)$, 
\item 对每个 $\delta > 0$, 
\begin{equation*}
\int_{|x| > \delta} |K(t,x)|\dd x \rightarrow 0 \quad (t \rightarrow 0).
\end{equation*}
\end{enumerate}
则

(a) 若 $f(x) \in L^{\infty}(\mathbb{R})$ 且在点 $x_{0}$ 连续, 则
\begin{equation*}
\int_{-\infty}^{\infty} K(t, x_{0} - y)f(y)\dd y \rightarrow f(x_{0}) \quad (t \rightarrow 0).
\end{equation*}

(b) 若 $f(x) \in L^{\infty}(\mathbb{R}) \cap C(\mathbb{R})$, 则在任意有限区间上一致地
\begin{equation*}
\int_{-\infty}^{\infty} K(t, x - y)f(y)\dd y \rightarrow f(x) \quad (t \rightarrow 0).
\end{equation*}

(c) 若 $1 \leq p < \infty$ 且 $f \in L^{p}(\mathbb{R})$, 则
\begin{equation*}
\left\| \int_{-\infty}^{\infty} K(t, x - y)f(y)\dd y - f(x) \right\|_{L^{p}} \rightarrow 0 \quad (t \rightarrow 0).
\end{equation*}

(d) 若 $1 \leq p \leq \infty$ 且 $f \in L^{p}(\mathbb{R})$, 则
\begin{equation*}
\left\| \int_{-\infty}^{\infty} K(t, x - y)f(y)\dd y \right\|_{L^{p}} \leq A\| f \|_{L^{p}}.
\end{equation*}
\end{theorem}

\begin{proof}
(a) 首先利用 (ii), 
\begin{align*}
\int_{-\infty}^{\infty} K(t, x_{0} - y)f(y)\dd y - f(x_{0})
& = \int_{-\infty}^{\infty} K(t, x_{0} - y)f(y)\dd y - f(x_{0})\int_{-\infty}^{\infty} K(t, x_{0} - y)\dd y \\
& = \int_{-\infty}^{\infty} K(t, x_{0} - y)\bigl(f(y) - f(x_{0})\bigr)\dd y.
\end{align*}
从而, 
\begin{align*}
& \left| \int_{-\infty}^{\infty} K(t, x_{0} - y)f(y)\dd y - f(x_{0}) \right| \\
& \qquad \leq \int_{|y - x_{0}| > \delta} |K(t, x_{0} - y)||f(y) - f(x_{0})|\dd y \\
& \qquad\quad + \int_{|y - x_{0}| \leq \delta} |K(t, x_{0} - y)||f(y) - f(x_{0})|\dd y \\
& \qquad \leq 2\| f \|_{\infty} \int_{|y - x_{0}| > \delta} |K(t, x_{0} - y)|\dd y \\
& \qquad\quad + \sup_{|x_{0} - y| \leq \delta} |f(y) - f(x_{0})| \int_{-\infty}^{\infty} |K(t, x_{0} - y)|\dd y.
\end{align*}
于是, 任取 $\varepsilon > 0$.因 $f(x)$ 在点 $x_0$ 连续, 故存在具有如下性质的 $\delta > 0$:
\begin{equation*}
|y - x_{0}| \leq \delta \quad \text{则} \quad |f(y) - f(x_{0})| < \varepsilon.
\end{equation*}
又由 (iii), 对此 $\delta$, 存在 $t_{0} > 0$ 使得
\begin{equation*}
0 < t < t_{0} \quad \text{则} \quad \int_{|y - x_{0}| > \delta} |K(t, x_{0} - y)|\dd y = \int_{|z| > \delta} |K(t, z)|\dd z < \varepsilon.
\end{equation*}
综合运用 (i) 可得, 对任意 $\varepsilon > 0$ 存在 $t_0 > 0$, 使得当 $0 < t < t_0$ 时
\begin{equation*}
\left| \int_{-\infty}^{\infty} K(t, x_{0} - y)f(y)\dd y - f(x_{0}) \right| \leq (2\| f \|_{\infty} + A)\varepsilon
\end{equation*}
故 (a) 得证.

(b) 令 $I = [-a, a]$, 则由假设 $f(x)$ 在 $I$ 上一致连续.由此与 (a) 的证法可得结论.

(c) 设 $f \in L^{p}(\mathbb{R})$.任取 $\varepsilon > 0$.利用 $C_{c}(\mathbb{R})$ 在 $L^{p}(\mathbb{R})$ 中稠密, 存在连续函数 $f_{0}$ 与正数 $R$ 使得
\begin{equation*}
\| f - f_{0} \|_{L^{p}} < \varepsilon
\end{equation*}
且 $\supp f_{0} \subset \{x;\ |x| \leq R\}$.利用 $L^{p}$ 的三角不等式
\begin{align*}
& \left\| \int_{-\infty}^{\infty} K(t, x - y)f(y)\dd y - f(x) \right\|_{L^{p}} \\
& \quad \leq \left\| \int_{-\infty}^{\infty} K(t, x - y)(f(y) - f_{0}(y))\dd y \right\|_{L^{p}} \\
& \quad\quad + \left\| \int_{-\infty}^{\infty} K(t, x - y)f_{0}(y)\dd y - f_{0}(x) \right\|_{L^{p}} + \| f_{0}(x) - f(x) \|_{L^{p}} \\
& \quad = I_{1} + I_{2} + I_{3}
\end{align*}
由 Young 不等式与 (i)
\begin{equation*}
I_{1} \leq \| K(t, \cdot) \|_{L^{1}} \| f - f_{0} \|_{L^{p}} \leq A\| f - f_{0} \|_{L^{p}}.
\end{equation*}
故
\begin{equation*}
I_{1} + I_{3} < (A + 1)\varepsilon.
\end{equation*}
关于 $I_{2}$, 利用 (b), 存在正数 $t_{0}$, 使得当 $0 < t < t_{0}$ 时, 对有界闭区间 $[-R-1, R+1]$ 上的任意 $x$, 
\begin{equation*}
\left| \int_{-\infty}^{\infty} K(t, x - y)f_{0}(y)\dd y - f_{0}(x) \right| < \frac{\varepsilon}{2(R+1)^{1/p}}.
\end{equation*}
从而
\begin{equation*}
\int_{|x| \leq R + 1}\left| \int_{-\infty}^{\infty} K(t, x - y)f_{0}(y)\dd y - f_{0}(x) \right|^{p}\dd x < \varepsilon^{p}.
\end{equation*}
其次, 利用 $f_{0}(x) = 0$($|x| \geq R$)与 Hölder 不等式,(i), 以及 Fubini 定理, 
\begin{align*}
J_{t} &:= \int_{|x| \geq R + 1}\left| \int_{-\infty}^{\infty} K(t, x - y)f_{0}(y)\dd y - f_{0}(x) \right|^{p}\dd x \\
&= \int_{|x| \geq R + 1}\left| \int_{-\infty}^{\infty} K(t, x - y)f_{0}(y)\dd y \right|^{p}\dd x \\
&\leq \int_{|x| \geq R + 1}\left(\int_{-\infty}^{\infty} |K(t, x - y)|\dd y\right)^{p/p'}\left(\int_{-\infty}^{\infty} |K(t, x - y)||f_{0}(y)|^{p}\dd y\right)\dd x \\
&\leq A^{p/p'}\int_{|x| \geq R + 1}\left(\int_{-\infty}^{\infty} |K(t, x - y)||f_{0}(y)|^{p}\dd y\right)\dd x \\
&= A^{p/p'}\int_{|x| \geq R + 1}\left(\int_{|y| \leq R} |K(t, x - y)||f_{0}(y)|^{p}\dd y\right)\dd x \\
&= A^{p/p'}\int_{|y| \leq R}\left(\int_{|x| \geq R + 1} |K(t, x - y)|\dd x\right)|f_{0}(y)|^{p}\dd y
\end{align*}
这里注意到当 $|y| \leq R$, $|x| \geq R+1$ 时 $|x-y| \geq |x| - |y| \geq (R+1) - R = 1$, 
\begin{align*}
J_{t} &\leq A^{p/p'}\int_{|y| \leq R}\left(\int_{|x - y| \geq 1} |K(t, x - y)|\dd x\right)|f_{0}(y)|^{p}\dd y \\
&= A^{p/p'}\int_{|y| \leq R}\left(\int_{|z| \geq 1} |K(t, z)|\dd z\right)|f_{0}(y)|^{p}\dd y \\
&= A^{p/p'}\int_{|y| \leq R} |f_{0}(y)|^{p}\dd y \int_{|z| \geq 1} |K(t, z)|\dd z \\
&= A^{p/p'}\| f_{0} \|_{L^{p}}^{p}\int_{|x| \geq 1} |K(t, x)|\dd x.
\end{align*}
由假设 (iii), 存在 $t_{1} > 0$ 使得
\begin{equation*}
J_{t} < \varepsilon^{p} \quad (0 < t < t_{1}).
\end{equation*}
故当 $0 < t < \min(t_{0}, t_{1})$ 时
\begin{equation*}
I_{2}^{p} < \varepsilon^{p} + \varepsilon^{p} < (2\varepsilon)^{p}.
\end{equation*}
从而
\begin{equation*}
I_{1} + I_{2} + I_{3} < (A + 3)\varepsilon
\end{equation*}
由此证得
\begin{equation*}
\lim_{t \rightarrow 0}\left\| \int_{-\infty}^{\infty} K(t, x - y)f(y)\dd y - f(x) \right\|_{L^{p}} = 0.
\end{equation*}

(d) 与 (c) 中 $I_{1}$ 的估计一样, 由 Young 不等式与 (i) 即得结论.
\end{proof}

\begin{example}\label{exa:1.1.10}
设 $\varphi(x) \in L^{1}(\mathbb{R})$, $\int_{-\infty}^{\infty} \varphi(x)\dd x = 1$, 并令 $\varphi_{t}(x) = t^{-1}\varphi(t^{-1}x)$.则 $K(t,x) = \varphi_{t}(x)$ 满足定理~\ref{thm:1.1.9} 的假设 (i),(ii),(iii).

事实上, (i) 由变量替换 $y = x/t$, 
\begin{equation*}
\int_{-\infty}^{\infty} |\varphi_{t}(x)|\dd x = \int_{-\infty}^{\infty}\left| \frac{1}{t}\varphi\left(\frac{x}{t}\right) \right|\dd x = \int_{-\infty}^{\infty} |\varphi(y)|\dd y.
\end{equation*}

(ii) 与 (i) 类似, 
\begin{equation*}
\int_{-\infty}^{\infty} \varphi_{t}(x)\dd x = \int_{-\infty}^{\infty} \frac{1}{t}\varphi\left(\frac{x}{t}\right)\dd x = \int_{-\infty}^{\infty} \varphi(y)\dd y = 1.
\end{equation*}

(iii) 因 $\varphi(x) \in L^{1}(\mathbb{R})$, 
\begin{equation*}
\int_{|x| > \delta} |\varphi_{t}(x)|\dd x = \int_{|x| > \delta}\left| \frac{1}{t}\varphi\left(\frac{x}{t}\right)\right|\dd x = \int_{|y| > \delta/t} |\varphi(y)|\dd y \rightarrow 0 \quad (t \rightarrow 0).
\end{equation*}
\end{example}
